\documentclass[11pt]{article}
\usepackage{amsmath,amssymb,amsthm,geometry,hyperref}
\geometry{margin=1in}
\newtheorem{theorem}{Theorem}
\newtheorem{lemma}[theorem]{Lemma}
\newtheorem{corollary}[theorem]{Corollary}
\newcommand{\E}{\mathbb E}
\newcommand{\KL}{\operatorname{KL}}
\newcommand{\relu}{\operatorname{ReLU}}
\title{A critical exponential tilt of a Gaussian ReLU neuron}
\date{18 November 2026; independent research draft}
\emergencystretch=3em
\author{ChatGPT-6 Astra\thanks{Independent AI research. Sole author. Verification scope and provenance are stated in the accompanying research collection.}}
\hypersetup{pdftitle={A critical exponential tilt of a Gaussian ReLU neuron},pdfauthor={ChatGPT-6 Astra}}
\begin{document}\maketitle
\begin{abstract}
The Gaussian ReLU parity statistic has a finite moment-generating
function at its positive critical boundary and a finite critical mean
in input dimension at least four. Above that mean, the scalar entropy
problem has a linear rate and no minimizing probability law. We prove
the exact local power-law tail under the critical tilt, use a cited
local large-deviation theorem, and derive a conditional-product limit with a
single exceptional neuron under scalar-statistic conditioning, together
with the leading evidence constant. A sampled conditioned network can
retain classification error approaching one quarter as the parity
dimension grows, despite an arbitrarily large enforced parity statistic.
Its posterior majority is correct. The joint conditional limit also
gives an independent posterior construction with total-variation error
tending to zero at fixed dimension. Its output fluctuations are stable
in dimension four, Gaussian with a logarithmic normalization in dimension
five, and Gaussian with the usual normalization in higher dimensions.
These results distinguish scalar
conditioning from conditioning on the individual labeled examples.
\end{abstract}
\noindent\textbf{Status.} The endpoint-integrability and entropy
nonattainment arguments are followed by a proved scalar-statistic
conditional condensation theorem, using the cited local big-jump result. The distinct posterior conditioned on all cube
margins is proved for even $r=d\ge6$ in the two later ReLU companions,
including its exact exceptional-neuron count. This manuscript is part of
the separate independent-research collection.

\section{The parity statistic and its critical tilt}
Let $x\in\{-1,1\}^d$, let $r\le d$ be even and at least two, and put
$\chi(x)=\prod_{i=1}^r x_i$. Write $\bar x=(1,x)$ and let
$U\sim N(0,I_{d+1})$ and $a\sim N(0,1)$ be independent. Define
\[
 G(U)=2^{-d}\sum_x\chi(x)\relu(U\cdot\bar x),\qquad Z=aG(U),
 \qquad \rho=\frac{\binom{r-2}{(r-2)/2}}{2^{r-1}\sqrt r}.
\]
The repository's sharp parity-correlation lemma establishes
$|G(U)|\le\rho\|U\|$, with equality precisely on the $2^r$ rays
\[
 v_s=(0,s_1/\sqrt r,\ldots,s_r/\sqrt r,0,\ldots,0),
 \qquad s\in\{-1,1\}^r.
\]
Both orientations of a line occur in this list. Let
$M(t)=\E e^{tZ}$ and $t_*=1/\rho$.

\begin{theorem}[Finite slope at the Gaussian boundary]\label{thm:boundary}
The moment generating function is finite for $|t|<t_*$ and infinite for
$|t|>t_*$. At the positive endpoint, for every integer $k\ge0$,
\[
 M^{(k)}(t_*-)<\infty\quad\Longleftrightarrow\quad d>2k+1.
\]
In particular, at $d=4$ the endpoint value and first derivative are finite,
while the second derivative is infinite. For every $d\ge4$ define
\[
 Q_*(dU\,da)=\frac{e^{t_*Z}}{M(t_*)}P_0(dU\,da),\qquad
 m_c=\E_{Q_*}Z=\frac{M'(t_*)}{M(t_*)}<\infty.
\]
The scalar Cram\'er rate $I(s)=\sup_t\{ts-\log M(t)\}$ has the exact
linear branch $I(s)=t_*s-\log M(t_*)$ for all $s\ge m_c$.
\end{theorem}

\begin{lemma}[A sharp corner at every maximizing ray]\label{lem:corner}
For $z\perp v_s$, define
\[
 h_s(z)=2^{-d-1}\sum_{x:\,\sum_{i\le r}s_ix_i=0}
                   |z\cdot\bar x|.
\]
This is a norm on the $d$-dimensional space $v_s^\perp$. For a sufficiently
small fixed $\eta>0$, whenever $R>0$ and $\|z\|\le\eta R$,
\[
 |G(Rv_s+z)|=\rho R-h_s(z).
\]
\end{lemma}
\begin{proof}
Since $r\ge2$, the affine part of $\relu(t)=(t+|t|)/2$ has zero
parity coefficient. At $v_s$, the inputs with zero preactivation have
parity $(-1)^{r/2}\chi(s)$, opposite to the sign of $G(v_s)$.
The nonzero preactivations keep their signs in a sufficiently small
conical neighborhood. Their contribution is exactly linear there.
Symmetry under permutations of relevant coordinates, sign changes of
irrelevant coordinates, and simultaneous input negation implies that its
gradient has no transverse component. Hence the zero preactivations
contribute precisely the displayed positive norm subtracted from the
absolute correlation.

To see positive definiteness, the zero-preactivation augmented input
vectors span $v_s^\perp$. Flipping an irrelevant coordinate spans its
coordinate axis. Averaging paired opposite core inputs supplies the bias
axis. Finally, differences of balanced subsets differing by one swap span
the relevant-coordinate hyperplane orthogonal to $(s_1,\ldots,s_r)$.
Thus $h_s(z)=0$ forces $z=0$ in $v_s^\perp$.
\end{proof}

\begin{proof}[Proof of Theorem~\ref{thm:boundary}]
Integrating the Gaussian readout gives
\[
 M(t)=(2\pi)^{-(d+1)/2}\int
       \exp\{-\tfrac12(\|U\|^2-t^2G(U)^2)\}\,dU.
\]
The sharp correlation bound proves finiteness inside the stated interval.
Outside it, continuity supplies an open cone on which the exponent grows
quadratically, proving divergence.

At $t=t_*$, outside small cones around the finitely many maximizing rays,
compactness of the unit sphere gives a strict quadratic gap. Within a
cone, use $U=Rv_s+z$ and Lemma~\ref{lem:corner}. Then
\[
 \|U\|^2-G(U)^2/\rho^2
 =2Rh_s(z)/\rho+\|z\|^2-h_s(z)^2/\rho^2.
\]
Norm equivalence and a smaller choice of $\eta$ bound this expression
above and below by positive constants times $R\|z\|$ on
$\|z\|\le\eta R$. After the substitution $z=w/R$, the integral with
an additional factor $G(U)^{2k}$ is comparable at infinity to
\[
 \int_1^\infty R^{2k-d}\,dR.
\]
For the lower comparison restrict $\|w\|$ to a fixed ball; for the upper
comparison integrate $e^{-c\|w\|}$ over all of $\mathbb R^d$.
This proves convergence exactly when $d>2k+1$.
Differentiating $e^{t^2G^2/2}$ $k$ times produces a polynomial in $G^2$
with nonnegative coefficients and positive leading coefficient at $t_*>0$,
so the same criterion applies to the endpoint derivative. For the final
claim, $\log M$ is convex and its largest finite left derivative is $m_c$;
therefore the Legendre supremum for $s\ge m_c$ is attained at $t_*$.
\end{proof}

\section{An entropy problem with no minimizing probability law}
Let $\Phi(U,a)=(a\relu(U\cdot\bar x))_x$ be the full one-neuron
output vector. All its components are integrable under $Q_*$ when $d\ge4$:
the ray-tube calculation controls $\E_{Q_*}(|a|\|U\|)$ by the same
integral as the first moment above. Cube symmetries then imply
\[
 \E_{Q_*}\Phi=m_c\chi.
\]
For clarity, a relevant input sign change can be compensated by changing
the readout by its parity character; this preserves both $P_0$ and $Z$.
Irrelevant input sign changes preserve the same laws without a readout
change. These symmetries force the mean vector into the one-dimensional
parity subspace.

\begin{theorem}[Nonattainment despite an exact small teacher]\label{thm:noattain}
Suppose $d\ge4$ and $s>m_c$. Each of the following entropy minimization
problems has value $t_*s-\log M(t_*)$ and has no minimizer:
\begin{align*}
 &\inf\{\KL(Q\|P_0):\E_Q Z\ge s\},\\
 &\inf\{\KL(Q\|P_0):\chi(x)\E_Q\Phi_x\ge s\text{ for every }x\},\\
 &\inf\{\KL(Q\|P_0):\E_Q\Phi=s\chi\}.
\end{align*}
The claims concern laws with finite indicated moments. The target is an
ordinary shallow ReLU parity function; $r=2$ already suffices if the input
has at least four coordinates.
\end{theorem}
\begin{proof}
The entropy inequality at $t_*$ gives
\[
 \KL(Q\|P_0)\ge t_*\E_Q Z-\log M(t_*).
\]
Equality requires $Q=Q_*$. Since its mean is $m_c<s$, no law can attain
the proposed lower bound in any of the three feasible sets.

It remains to match that bound. Let $\kappa=G(v_{\boldsymbol1})$.
For $R>0$, define $\Pi_R$ as the uniform mixture over $s\in\{-1,1\}^r$
of Gaussians with identity covariance and respective means
\[
 U=\sqrt{R/\rho}\,v_s,\qquad
 a=\operatorname{sgn}(\kappa)\chi(s)\sqrt{R/\rho}.
\]
Convexity of relative entropy and the Gaussian shift formula give
$\KL(\Pi_R\|P_0)\le R/\rho$.
Without the Gaussian smoothing, the mean output is exactly $R\chi$,
by Walsh projection over the orientations. ReLU's Lipschitz property
and independence of the readout imply that smoothing changes each output
by at most $\sqrt{R/\rho}\sqrt{d+1}$. The same symmetries as above force
the smoothed mean to be $m_R\chi$, where $m_R=R+O(\sqrt R)$.

For large $R$, put $\eta_R=(s-m_c)/(m_R-m_c)\in(0,1)$ and
$Q_R=(1-\eta_R)Q_*+\eta_R\Pi_R$. Then $\E_{Q_R}\Phi=s\chi$ exactly,
and
\[
 \limsup_{R\to\infty}\KL(Q_R\|P_0)
 \le\KL(Q_*\|P_0)+(s-m_c)/\rho=t_*s-\log M(t_*).
\]
This proves all three infima. Notice also that $Q_R\to Q_*$ in total
variation while their mean outputs remain $s\chi$: a vanishing amount of
probability at increasingly large weights carries the missing mean.
\end{proof}

The full-output Cram\'er rate consequently obeys
$I_{\Phi}(s\chi)=t_*s-\log M(t_*)$ for $s\ge m_c$.
The lower bound uses the parity dual direction; the upper bound follows
by applying the entropy inequality for every exponential tilt to $Q_R$.
This does not invalidate a Cram\'er large-deviation principle. It prevents
identifying the conditioned limit with an ordinary entropy-minimizing
one-neuron probability law beyond the critical mean.

\section{The local tail and a single exceptional neuron}
The scalar-statistic conditioning problem admits a complete description.
Its general probabilistic input is the local large-deviation theorem of
Denisov--Dieker--Shneer. The associated conditional-product phenomenon
was established by Armend\'ariz--Loulakis. We verify the local
large-deviation hypotheses for the critical Gaussian ReLU law and give
a direct conditional-density argument including the overshoot.
This is not a new general
one-big-jump theorem, and the conditioning is not a collection of
individual labels.

Put $\alpha=(d-1)/2>1$ and
\[
 C=\frac{(2\pi)^{-(d+1)/2}}{M(t_*)}
 \frac{\rho^\alpha}{d-1}
 \sum_s\int_{v_s^\perp}e^{-h_s(t)/\rho}\,dt.
 \tag{10}
\]
This constant is positive and finite by the transverse norm property.

\begin{lemma}[Density and marks in the critical tail]\label{lem:localtail}
On $(0,\infty)$ the law of $Z$ under $Q_*$ has a density $f_*$ with
\[
 f_*(z)\sim\alpha C z^{-\alpha-1},\qquad
 Q_*(Z>z)\sim C z^{-\alpha}.
 \tag{11}
\]
Its negative tail is exponentially small. Conditional on $Z=z$ as
$z\to\infty$, the parameter vector approaches one of the $2^r$
critical rays, uniformly among them, and
\[
          \frac{a^2+\|U\|^2}{z}\longrightarrow\frac2\rho
 \quad\text{in probability}.                          \tag{12}
\]
A possible atom at $Z=0$ is harmless and is not excluded.
\end{lemma}
\begin{proof}
Inside a small critical cone write
\[
 a=\sigma_s A,\quad U=Rv_s+u,\quad c=AR,\quad v=A-R,
 \quad t=\sqrt c\,u,
\]
where $A,R>0$, $u\perp v_s$, and $\sigma_sG(v_s)=\rho$.
Let $q=A/\sqrt c=\sqrt{A/R}$. The exact corner identity gives
\[
 Z=\rho c-qh_s(t),\qquad
 D_0:=\tfrac12(a^2+\|U\|^2)-Z/\rho
       =\tfrac12v^2+\frac{\|t\|^2}{2c}+\frac q\rho h_s(t),
\]
with Jacobian $c^{-d/2}/(A+R)$ for $dc\,dv\,dt$.
On this cone $q$ is bounded above and below and
$D_0\ge c_1(v^2+\|t\|)$.
Furthermore
\[
 \frac{\partial q}{\partial c}
       =-\frac{qv}{2c\sqrt{v^2+4c}},\qquad
 \frac{\partial Z}{\partial c}=\rho-h_s(t)\frac{\partial q}{\partial c}.
\]
Shrinking the incoming cone ensures $qh_s(t)\le\eta c$ with
$\eta<\rho$. Thus $\partial Z/\partial c\ge\rho-\eta/2>\rho/2$,
and on $Z=z>0$ one has $z/\rho\le c\le z/(\rho-\eta)$.
The change from $c$ to $z$ is therefore regular.

For fixed $(v,t)$ and $z\to\infty$, one has
$c\sim z/\rho$, $q\to1$, $(A+R)/\sqrt c\to2$, and
$\partial Z/\partial c\to\rho$. After multiplying the resulting
density by $z^{\alpha+1}$ its integrand converges to
\[
 \frac{(2\pi)^{-(d+2)/2}}{M(t_*)}
       \frac{\rho^\alpha}{2}
                      e^{-v^2/2-h_s(t)/\rho}.
\]
It is dominated on the whole cone by a constant times
$e^{-c_1(v^2+\|t\|)}$. The cone domains exhaust the transverse
space and the real $v$ axis. Dominated convergence, followed by the
Gaussian integral in $v$, gives exactly $\alpha C$ in (11).

For clarity, the part outside these finitely many cones is negligible
also at the density level. The homogeneous quadratic exponent $D_0$
has a positive gap there on the unit sphere. Integrate the readout at
$a=z/G(U)$; the density integrand is $|G(U)|^{-1}e^{-D_0}$ times
its normalizing constant. The gap bounds it by
\[
 e^{-c_2z}\,|G(U)|^{-1}
           e^{-c_3\{\|U\|^2+z^2/G(U)^2\}}.
\]
Here $a^2+\|U\|^2\ge2z/\rho$. The inequality
$y^{-1}e^{-c_3z^2/y^2}\le C_3/z$ for $z>0$, $y>0$ bounds the
remaining integral by $C_4/z$. Points with $G(U)=0$ contribute only
to the atom at zero. This proves the asserted density asymptotic.
Integration gives the tail. Readout symmetry before tilting implies
$f_*(-z)=e^{-2t_*z}f_*(z)$ on $z>0$, proving the negative-tail claim.
The same dominated calculation makes $(v,t)$ tight given large $z$,
which yields (12) and convergence to the critical rays. Relevant input
sign changes, compensated in the readout, preserve $Z$ and act
transitively on those rays; their limiting probabilities are equal.
\end{proof}

For example, at $d=r=4$, orthonormal transverse coordinates give
\[
 h_s(t)/\rho=\sum_{i=1}^3\max(|b|/2,|t_i|),\qquad
 \int e^{-h_s(t)/\rho}dt
       =32\int_0^\infty(1+u)^3e^{-3u}du=832/27.
\]
The earlier importance estimates $M(8)\approx1.178$ and
$m_c\approx0.1805$ give $C\approx0.0623$; these displayed decimal
values are estimates, not interval certificates.

\begin{theorem}[Scalar conditioning above the critical mean]\label{thm:scalar}
Fix $d\ge4$, an even $2\le r\le d$, and $s>m_c$. Let
$\delta=s-m_c$ and $S_N=\sum_{i=1}^N Z_i$ for independent prior
neurons. Then
\[
 P_0^N(S_N\ge Ns)
 \sim e^{-NI(s)}\frac{\alpha C}{t_*\delta^{\alpha+1}}N^{-\alpha}.
 \tag{13}
\]
Under this conditional law, the overshoot $S_N-Ns$ converges in
law to an exponential variable of rate $t_*$. After moving the neuron
with largest $Z_i$ to the last position, the joint law of all other
$N-1$ parameter vectors converges in total variation to $Q_*^{N-1}$.
For the largest neuron $J_N$,
\[
 Z_{J_N}/N\to\delta,\qquad
 (a_{J_N}^2+\|U_{J_N}\|^2)/N\to2\delta/\rho,
 \qquad \max_{i\ne J_N}\|\theta_i\|/\sqrt N\to0
 \tag{14}
\]
in conditional probability. Its ray is asymptotically uniform among
the $2^r$ critical rays. In particular exactly one neuron exceeds
$\epsilon\sqrt N$ in norm with probability tending to one whenever
$0<\epsilon<\sqrt{2\delta/\rho}$.
\end{theorem}
\begin{proof}
Apply the cited local large-deviation result to the centered critical variable
$X=Z-m_c$. Lemma~\ref{lem:localtail} gives a regularly varying positive
local density and an exponentially bounded negative tail. Choose
$1<\kappa<\min\{\alpha,2\}$. Then $\E_{Q_*}|X|^\kappa<\infty$;
for each fixed interval length $h>0$ its local probability is almost
decreasing, regularly varying, and insensitive to shifts of order
$x^{1/\kappa}=o(x)$. These verify Corollary 2.1 of
\href{https://arxiv.org/abs/math/0703265}{Denisov--Dieker--Shneer (2008)}.
It supplies, uniformly for $y\ge N\delta$,
\[
 Q_*^N\{S_N-Nm_c\in(y,y+h]\}
       \sim Nh f_*(y).
 \tag{15}
\]
The constant mean shift in the one-step density does not alter the
asymptotic. The product-bulk conclusion is consistent with Theorem 1 of
\href{https://arxiv.org/abs/0912.1516}{Armend\'ariz--Loulakis (2011)};
the direct density argument below additionally tracks the overshoot.

Here are details of the further exponential reweighting, which is
necessary because the desired conditional law is under $P_0$, not $Q_*$.
Changing measure gives the residual factor
\[
 D_N=\E_{Q_*^N}\left[e^{-t_*(S_N-Ns)}\mathbf1_{\{S_N\ge Ns\}}\right],
 \qquad P_0^N(S_N\ge Ns)=e^{-NI(s)}D_N.
\]
Partition the nonnegative overshoot axis into intervals of length $h$.
For each bounded overshoot (15) gives probability asymptotic to
$Nh f_*(N\delta)$. The same uniform estimate and regular variation
bound every interval probability by $C_hNh f_*(N\delta)$ for all
large $N$. Thus exponential weights give a summable dominating sequence.
Upper and lower endpoint Riemann sums, with $N\to\infty$ first and
$h\downarrow0$ afterward, yield
\[
 D_N\sim Nf_*(N\delta)\int_0^\infty e^{-t_*u}\,du
                 =Nf_*(N\delta)/t_*.
\]
This proves (13). Starting the same sums at a fixed overshoot $b\ge0$
gives limiting survival probability $e^{-t_*b}$.

Here is the direct product argument. Let $\mu$ be the law of $Z$
under $Q_*$. Given a proposed bulk $z=(z_1,\ldots,z_{N-1})$, put
$b(z)=Ns-\sum_i z_i$ and $m(z)=\max_i z_i$.
After swapping the maximum to the last position, the joint conditional
law of this scalar bulk and the overshoot $u\ge0$ has density
\[
 \frac{N}{D_N}e^{-t_*u}f_*(b(z)+u)
                  \mathbf1_{\{b(z)+u>m(z)\}}             \tag{15a}
\]
with respect to $\mu^{N-1}(dz)\,du$. The maximum is positive and
unique almost surely on the conditioning event, so exchangeability
gives exactly the factor $N$, even if $\mu$ has an atom at zero.
Under the reference probability law
$\mu^{N-1}\otimes\operatorname{Exp}(t_*)$, one has
$b(z)/N\to\delta$, $m(z)/N\to0$, and $u/N\to0$ in probability.
The finite first moment gives the first two claims. Regular variation
of $f_*$ and $D_N\sim Nf_*(N\delta)/t_*$ show that the density ratio
in (15a) relative to this reference law tends to one in probability.
Its integral is one. For any such nonnegative density ratio $R_N$,
$\E|R_N-1|=2\E(1-R_N)_+\to0$ by bounded convergence in probability.
This proves the stronger joint total-variation convergence of bulk and
overshoot to the indicated independent product.
The conditional parameter law given each $Z_i$ is identical under
$P_0$ and $Q_*$; applying these independent mark kernels cannot increase
total variation, so the same conclusion holds for the parameter vectors.

The finite first moment gives the bulk law of large numbers. Combined
with tight overshoot, it implies $Z_{J_N}/N\to s-m_c$.
Lemma~\ref{lem:localtail} gives the marked large-neuron limits in (14).
The critical parameter second moment is finite for $d>3$, by the same
ray-coordinate integral as the first moment of $Z$. Thus
$NQ_*(\|\theta\|>\epsilon\sqrt N)\to0$ for each fixed $\epsilon>0$.
A union bound under the product bulk law, and its total-variation
approximation, prove the final limit in (14).
\end{proof}

The theorem completes the scalar-conditioning application left open in
the first-round draft. Individual-label conditioning is different: full
cube margins require compatible contributions in all antipodal directions,
which is why the later full-output companions have exponentially many
exceptional neurons rather than one. None of these scalar-conditioning
statements asserts an OOD mistake lower bound.

\section{A large parity statistic need not classify the parity}
The scalar-conditioned network has an explicit limiting output, which
also makes the limitation of this conditioning visible. Write
$\phi_v(x)=\relu(v\cdot\bar x)$ and
$\sigma_s=\operatorname{sgn}G(v_s)$. For the original normalization
$F_N=N^{-1}\sum_i\Phi(\theta_i)$, the preceding theorem gives
\[
 F_N\ \Longrightarrow\ m_c\chi+
             \frac{\delta}{\rho}\sigma_S\phi_{v_S}
 \quad\text{in }\mathbb R^{\{-1,1\}^d},                 \tag{16}
\]
where $S$ is uniform among the $2^r$ critical orientations. The first
term follows from the finite-mean product bulk law; homogeneity and
(12) give the second. Thus one posterior draw has a random limiting
function even though the ensemble is invariant under all source-label
symmetries.

\begin{corollary}[Nonvanishing classification error of a scalar-conditioned draw]
Let $r\ge4$ be even and $d\ge r$. If
\[
                 \delta>\frac{m_c\rho\sqrt r}{4},       \tag{17}
\]
the uniform-input classification error of one scalar-conditioned
posterior draw converges in probability to
\[
 q_r=\frac14-\frac1{2^{r+1}}\binom r{r/2}.
 \tag{18}
\]
In particular $q_r\to1/4$ as $r\to\infty$, after the width limit.
For every fixed input, the posterior probability that the draw
misclassifies it also converges to $q_r$. Consequently posterior
majority classification is eventually correct at every input, while a
sampled network retains error (18). Its posterior mean output converges
to $s\chi$.
\end{corollary}
\begin{proof}
Put $r=2m$ and $y_i=S_ix_i$ for $i\le r$. The signed margin in (16)
is
\[
 m_c+\frac\delta\rho(-1)^{m-1}\chi(y)
             \relu\left(\frac{\sum_i y_i}{\sqrt r}\right).
\]
Here $m_c>0$: the untilted statistic is symmetric and nondegenerate,
so its finite positive-endpoint tilted mean is strictly positive.
If $k$ coordinates of $y$ are negative, the second term can be negative
only when $k<m$ and $k\equiv m\pmod2$. In that case $m-k\ge2$,
so its absolute value is at least $4\delta/(\rho\sqrt r)>m_c$.
At every other input the margin is at least $m_c$. All limiting margins
are therefore bounded away from zero. Exactly
\[
 \sum_{\substack{k<m\\k\equiv m\ (2)}}\binom r k
       =2^{r-2}-\frac12\binom r m
\]
inputs in the relevant cube have a negative margin; the identity follows
by pairing $k$ with $r-k$ in the parity-restricted binomial sum.
This count is independent of $S$, so (16) proves convergence of the
sampled classification error. For fixed $x$, uniform $S$ makes $y$
uniform, proving the posterior-probability assertion and, since
$q_r<1/2$, the majority assertion. The finite cube permits the latter
conclusion simultaneously at all inputs for sufficiently large widths.

Finally the exact conditional symmetries force
$\E[F_N\mid S_N\ge Ns]=c_N\chi$, with
$c_N=\E[S_N/N\mid S_N\ge Ns]$.
The uniform interval bound in the proof of (13), summed with exponential
weights, bounds the conditional overshoot tail by $C e^{-t_*b}$
uniformly for all large $N$ and $b\ge0$. Its means are bounded, hence
$c_N\to s$.
\end{proof}

This is a failure of learning from only a global parity statistic, not
a failure of the ordinary posterior conditioned on labeled examples.
It also separates sampling a network from taking the posterior majority;
the latter has no classification failure in this example.

\section{An asymptotically exact independent construction}
The preceding density calculation supplies more than convergence of the
exceptional neuron's direction. Fix orthonormal transverse coordinates
at each ray, related by the input-sign symmetries, and let
\[
 H_r=\int_{v_s^\perp}e^{-h_s(t)/\rho}\,dt.
\]
The conditional law of its coordinates $(s,v,t)$ given $Z=z$ converges
in total variation, as $z\to\infty$, to
\[
 \mathcal M(ds\,dv\,dt)
 =\frac1{2^r}\frac{e^{-v^2/2}}{\sqrt{2\pi}}\,dv\,
                         H_r^{-1}e^{-h_s(t)/\rho}\,dt.    \tag{19}
\]
Assign a cemetery mark outside the chosen critical cones. Its conditional
probability tends to zero. Inside the cones the proof of
Lemma~\ref{lem:localtail} gives convergence of the normalized densities
with an integrable dominating function. This proves (19), including the
vanishing cemetery mass, by $L^1$ convergence of densities.

\begin{theorem}[Independent scalar-posterior construction]
Fix even $r\ge4$, $d\ge r$, and $s>m_c$. The following distribution
on width-$N$ networks has total-variation distance tending to zero from
$P_0^N(\,\cdot\mid S_N\ge Ns)$.
\begin{enumerate}
\item Draw $N-1$ independent parameter vectors from $Q_*$, independently
draw $u\sim\operatorname{Exp}(t_*)$ and a mark $(s_0,v,t)\sim\mathcal M$,
and set $z=Ns-\sum_{i<N}Z_i+u$.
\item If $z\le0$, restart. Otherwise solve for $c>0$ in
\[
 z=\rho c-\frac{A(c,v)}{\sqrt c}h_{s_0}(t),\qquad
 A(c,v)=\frac{\sqrt{v^2+4c}+v}{2},\quad
 R(c,v)=\frac{\sqrt{v^2+4c}-v}{2}.                       \tag{20}
\]
Construct the last neuron as
$a=\sigma_{s_0}A(c,v)$ and
$U=R(c,v)v_{s_0}+t/\sqrt c$.
\item Restart if this neuron lies outside its prescribed critical cone
or if $z\le\max_{i<N}Z_i$. Otherwise swap it with a uniformly selected
index among $1,\ldots,N$.
\end{enumerate}
The probability of a restart tends to zero. Before the final swap,
the output statistic is exactly $Ns+u$ on every accepted draw.
For fixed $r,d$, the idealized construction uses an expected $O_{r,d}(N)$
independent elementary proposals and scalar-root operations as $N\to\infty$.
There is no asserted dimension-uniform cost, finite-$N$ variation rate,
or bound on the bit complexity of real-arithmetic operations.
\end{theorem}
\begin{proof}
First, (20) has a unique positive solution for every $z>0$.
Put $q=A/\sqrt c$ and $H=h_{s_0}(t)\ge0$. The function
$Z(c)=\rho c-qH$ tends to infinity as $c\to\infty$ and has
nonpositive upper limit as $c\downarrow0$. Moreover
$|cq'/q|\le1/2$, by the derivative already computed. Whenever
$Z(c)>0$, it follows that
\[
 Z'(c)\ge\rho-\frac{qH}{2c}
             =\frac\rho2+\frac{Z(c)}{2c}>\frac\rho2.
\]
Every positive level is therefore crossed exactly once.

Under the proposed independent bulk and overshoot law, $z/N\to\delta$
and $\max_{i<N}Z_i/N\to0$ in probability. Marks from the fixed law
$\mathcal M$ are tight. Thus $c/N\to\delta/\rho$, $A/R\to1$,
and $\|t\|/(\sqrt c R)\to0$ in probability. These limits put the
constructed neuron inside any of the fixed small critical cones with
probability tending to one, and make $z$ the largest statistic.
The exact corner identity then verifies $Z=z$. Restarting changes
the candidate law by total-variation distance tending to zero.

Equation (15a) proves joint total-variation convergence of the true
swapped bulk and overshoot to the proposed independent product, also
after lifting the bulk scalars to their conditional parameter laws.
Under either of these laws the exceptional statistic tends to infinity
in probability. Equation (19) therefore allows its exact conditional
mark kernel to be replaced by $\mathcal M$, again with vanishing
total-variation error. Applying the deterministic reconstruction (20)
cannot increase this distance. Finally exchangeability makes the true
largest index uniform and independent of its swapped configuration;
the last random swap recovers the labeled network law.

For completeness, the required independent distributions can be
sampled at a finite fixed-dimension expected cost without using the
conditioned-network Markov chain. The exact critical-law rejection
envelope is given in the companion on exponential ReLU condensation;
its acceptance probability is $(r-2)M(t_*)/r\ge(r-2)/r$.
The same construction permits $d-r$ irrelevant Gaussian coordinates.
Because $h_{s_0}$ is a norm in a fixed transverse space, choose
$a_0>0$ with $h_{s_0}(t)/\rho\ge a_0\|t\|_1$ in the chosen
orthonormal coordinates. A product Laplace proposal with density
proportional to $e^{-a_0\|t\|_1}$, accepted with probability
$e^{-h_{s_0}(t)/\rho+a_0\|t\|_1}$, samples its mark density exactly.
Its acceptance probability is a fixed positive constant. The normal
and exponential marks are direct draws, and only one root is required
per whole-network proposal. Since the overall acceptance tends to one,
the stated fixed-dimension expected proposal count follows.
\end{proof}

The construction does not use the unknown evidence or critical
normalization in its acceptance tests. It is asymptotically exact, not
an exact finite-width posterior sampler. Its conditioning contains one
scalar statistic; the many-label posterior has a different exceptional
configuration and is not sampled by this procedure.

\section{Dimension-dependent output fluctuations}
The local tail also identifies why convergence at dimension four can be
slow. Gaussian and stable condensate fluctuations are established
phenomena in heavy-tailed conditioned sums; see, for example,
\href{https://arxiv.org/abs/0912.1793}{Armend\'ariz--Grosskinsky--Loulakis},
Section 2.3, including its discussion of infinite-variance bulk laws.
The result here computes the dimension threshold, spectral measure, and
output projection for the Gaussian ReLU statistic. It does not claim
a new general stable-limit principle. Let $f_s=\sigma_s\phi_{v_s}$, so
$\langle\chi,f_s\rangle_{L^2(X)}=\rho$, and define the projection
\[
 \mathcal P_s y=y-\rho^{-1}f_s\langle\chi,y\rangle_{L^2(X)}.
\]
This need not be an orthogonal projection. It removes the one scalar
component constrained by the conditioning event.

\begin{theorem}[Stable and Gaussian fluctuation regimes]
Keep fixed $d\ge4$ and even $2\le r\le d$, and $s>m_c$.
Let $S_N^{\rm ray}$ be the nearest critical orientation of the largest
neuron, assigning any value off the event where it is well defined.
Put $\alpha=(d-1)/2$ and
\[
 b_N=\begin{cases}
 N^{2/3},&d=4,\\
 \sqrt{N\log N},&d=5,\\
 \sqrt N,&d\ge6.
 \end{cases}
\]
Under scalar conditioning, the following joint convergence holds:
\[
 \left(S_N^{\rm ray},\frac N{b_N}
    \left[F_N-m_c\chi-\frac\delta\rho f_{S_N^{\rm ray}}\right]\right)
       \Longrightarrow (S,\mathcal P_S Y),                \tag{21}
\]
where $S$ is uniform among the critical orientations and independent
of $Y$. The distribution of $Y$ is as follows.
\begin{enumerate}
\item For $d=4$, it is the centered $3/2$-stable vector with
characteristic exponent
\[
 \log\E e^{i\xi\cdot Y}
 =\frac{\alpha C}{2^r}\sum_s\int_0^\infty
  \left(e^{it\xi\cdot f_s/\rho}-1-it\xi\cdot f_s/\rho\right)
                           t^{-\alpha-1}\,dt.            \tag{22}
\]
Here $\xi\cdot y$ denotes the ordinary Euclidean pairing on the
finite output-vector space.
\item For $d=5$, it is centered Gaussian with covariance
\[
                 \Gamma=\frac{C}{2^r\rho^2}
                                   \sum_s f_sf_s^\top.   \tag{23}
\]
\item For $d\ge6$, it is centered Gaussian with covariance
$\Sigma=\operatorname{Cov}_{Q_*}(\Phi)$, which is finite.
\end{enumerate}
The stable or Gaussian vector can be degenerate on some output
directions. No finite-width error bound is asserted.
\end{theorem}
\begin{proof}
We first identify the independent critical-bulk limit. The density and
mark lemma implies multivariate regular variation of $\Phi$ under $Q_*$:
on sets bounded away from zero its tail measure is
\[
 \nu(B)=\frac{\alpha C}{2^r}\sum_s\int_0^\infty
                    \mathbf1_{\{tf_s/\rho\in B\}}
                                      t^{-\alpha-1}\,dt. \tag{24}
\]
To justify that no additional output tails are omitted, write
$D_0=\|\theta\|^2/2-Z/\rho$. Outside any fixed conical neighborhood
of the critical rays, $D_0\ge c\|\theta\|^2$ by compactness of the
unit sphere. The same gap holds on
$Z\le\rho\|\theta\|^2/4$. Since $\|\Phi\|\le C_d\|\theta\|^2$,
large output on either of these sets has an exponentially small tail.
On the remaining ray tubes, the local coordinate calculation gives
$\Phi/Z\to f_s/\rho$ with the stated uniform ray weights. It also
gives the bound $Q_*(\|\Phi\|>t)\le C't^{-\alpha}$ for large $t$.
These facts prove (24) and the required truncated-moment bounds.

For $d=4$, apply (24) to
$e^{i\xi\cdot y}-1-i\xi\cdot y$ with scale $b_N=N^{1/\alpha}$.
The contribution from $\|y\|\le\eta$ is bounded by $O(\eta^{2-\alpha})$
using the truncated second moment. The contribution from
$\|y\|>R$ is bounded by $O(R^{1-\alpha})$ using the tail first moment.
On the intervening annulus, tail-measure convergence applies directly.
Thus the characteristic exponent of
$b_N^{-1}\sum_{j<N}(\Phi_j-m_c\chi)$ tends to (22).
The quadratic error in expanding the logarithm is negligible because
$N/b_N^2\to0$ and the first moment is finite. This proves the stable
bulk limit without assuming a finite variance.

For $d=5$, (24) with $\alpha=2$ instead gives
\[
 \E[\Phi\Phi^\top\mathbf1_{\{\|\Phi\|\le R\}}]
       \sim \frac{2C\log R}{2^r\rho^2}\sum_s f_sf_s^\top
\]
in the sense of entrywise leading logarithmic coefficients; bounded
centering changes no coefficient. With $b_N=\sqrt{N\log N}$,
$\log b_N/\log N\to1/2$, giving (23).
The chance of a summand above $b_N$ is $O(1/\log N)$ after a union
bound. The centered mean removed by this truncation, divided by $b_N$,
is also $O(1/\log N)$. For each fixed $\eta>0$, the variance of
summands between $\eta b_N$ and $b_N$, divided by $b_N^2/N$, tends
to zero. The elementary triangular-array Lindeberg theorem therefore
gives the claimed Gaussian bulk limit. For $d\ge6$, the ordinary
finite-variance multivariate central limit theorem applies directly.

Finally the joint product law proved from (15a) makes these bulk limits
valid under the conditioned law, independently of the exceptional
orientation. Its tight mark law gives
\[
 \Phi(\theta_{J_N})=\rho^{-1}Z_{J_N}f_{S_N^{\rm ray}}+O_P(1).
\]
Indeed $\Phi=c f_s+O_P(1)$ and $Z=\rho c+O_P(1)$ in the tight
$(v,t)$ coordinates. Since
$Z_{J_N}=Ns+u-\sum_{i\ne J_N}Z_i$ with tight $u$, substitution
into the total output gives exactly $\mathcal P_{S_N^{\rm ray}}$
applied to the centered bulk sum, up to $O_P(1)$.
Dividing by $b_N\to\infty$ proves (21).
\end{proof}

In particular dimension-four output fluctuations are of order
$N^{-1/3}$, rather than $N^{-1/2}$. The logarithmic boundary at
dimension five and the finite-variance regime above it concern the
critical ordinary neurons. They do not follow from a finite-variance
approximation imposed on the exceptional neuron.

\section{An angular-volume criterion for the critical tail}
The dimension thresholds above come from the geometry of the maximizing
correlation directions. The following elementary criterion separates that
geometry from the particular parity calculation. It is a Laplace-method
calculation, not a new general theory of heavy-tailed condensation.

Let $U\sim N(0,I_{d+1})$, $a\sim N(0,1)$ independently, and now let
$G:\mathbb R^{d+1}\to\mathbb R$ be any continuous positively
homogeneous function of degree one. Write
$g(\omega)=G(\omega)$ for $\omega\in S^d$ and assume
$\rho=\max_{S^d}|g|>0$. Surface measure below is unnormalized.
Define the angular near-maximum volume
\[
 V(u)=\operatorname{area}\left\{\omega\in S^d:
        |g(\omega)|>0,\quad |g(\omega)|^{-1}-\rho^{-1}\le u\right\}.
\]
Suppose, for some $L>0$ and $\kappa>0$,
\[
                         V(u)\sim L u^\kappa\qquad(u\downarrow0).
 \tag{25}
\]
The zero set of $g$ is allowed to have positive surface measure.

\begin{theorem}[Angular volume determines endpoint moments]
Put $Z=aG(U)$ and $t_*=1/\rho$, and set
$\alpha=\kappa-(d+1)/2$. For each integer $k\ge0$,
\[
        M^{(k)}(t_*-)<\infty\quad\Longleftrightarrow\quad\alpha>k.
 \tag{26}
\]
If $\alpha>0$, the critical law $Q_*=e^{t_*Z}P_0/M(t_*)$ exists,
and its density on the positive half-line satisfies
\[
 f_*(z)\sim A z^{-\alpha-1},\qquad
 A=\frac{(2\pi)^{-(d+2)/2}}{M(t_*)}
       \sqrt{\frac\pi2}\,
       \rho^{-(d+1)/2}L\Gamma(\kappa+1).
 \tag{27}
\]
Consequently $Q_*(Z>z)\sim (A/\alpha)z^{-\alpha}$.
If $\alpha>1$, its mean $m_c$ is finite and the scalar entropy
problem above $m_c$ has the same linear, unattained infimum as in
Theorem~\ref{thm:noattain}. This last assertion concerns only the scalar
constraint, without an additional symmetry hypothesis on output vectors.
\end{theorem}
\begin{proof}
After integrating $a$, polar coordinates express the endpoint integral,
up to a positive constant, as
\[
 \int_0^\infty R^d\int_{S^d}
       \exp\{-R^2[1-g(\omega)^2/\rho^2]/2\}\,d\omega\,dR.
\]
Near the maximum,
$1-g^2/\rho^2\sim2\rho(|g|^{-1}-\rho^{-1})$.
Integration by parts in the angular sublevel volume (25) gives an
angular integral asymptotic to a positive constant times $R^{-2\kappa}$.
The part away from the maximum is exponentially small. Inserting
$G(U)^{2k}=R^{2k}g(\omega)^{2k}$ changes this asymptotic by a
positive constant times $R^{2k}$. The radial tail therefore converges
exactly when $d+2k-2\kappa<-1$. As in the proof of
Theorem~\ref{thm:boundary}, differentiating after integrating $a$
produces a polynomial in $G^2$ with positive leading coefficient and
nonnegative other coefficients at the positive endpoint. This proves
(26), including logarithmic divergence at equality.

For the density, integrate $a=z/G(U)$, and then use polar coordinates.
With $q=|g(\omega)|>0$ this gives
\[
 f_*(z)=\frac{(2\pi)^{-(d+2)/2}}{M(t_*)}e^{z/\rho}
 \int_{S^d}\frac1q\int_0^\infty
 R^{d-1}\exp\left\{-\frac{R^2}2-\frac{z^2}{2q^2R^2}\right\}
                                           dR\,d\omega. \tag{28}
\]
Directions with $q=0$ contribute only to the atom at zero.
The elementary radial saddle estimate is
\[
 \int_0^\infty R^{d-1}e^{-R^2/2-z^2/(2q^2R^2)}dR
 \sim\sqrt{\frac\pi2}\,(z/q)^{(d-1)/2}e^{-z/q}.
 \tag{29}
\]
It is uniform for $q$ in a fixed neighborhood of $\rho$.
For example set $R=\sqrt{z/q}\,y$; the exponent has its unique
minimum at $y=1$, with value $z/q$ and second derivative $4z/q$.
Taylor expansion on a shrinking neighborhood of one, followed by
Gaussian integration, gives (29). The complement has a strict
exponential gap. This proof also gives an upper bound by a constant
multiple of its right side whenever $z/q\ge1$.

The angular region $q\le\rho-\epsilon$ contributes exponentially
less than the main term: the last bound, with its fixed power of
$q^{-1}$, is dominated by $C e^{-c_\epsilon z}$ after factoring
out $e^{-z/\rho}$. In the remaining region the factor
$q^{-(d+1)/2}$ tends to $\rho^{-(d+1)/2}$ at the maximizing set.
Finally (25) and integration by parts give
\[
 \int_{S^d}e^{-z(q^{-1}-\rho^{-1})}d\omega
       \sim L\Gamma(\kappa+1)z^{-\kappa}.
\]
Substitution in (28)--(29) proves (27). Integrating the density gives
the tail. Readout symmetry again makes the negative density
$e^{-2t_*z}f_*(z)$.

For the entropy assertion, its lower bound and nonattainment follow
from the critical exponential-tilt inequality exactly as before.
For a matching sequence let $B_R=\{R<Z\le R+1\}$ and
$\Pi_R=Q_*(\,\cdot\mid B_R)$. Equation (27) gives
$Q_*(B_R)\sim A R^{-\alpha-1}$, whence
\[
 \operatorname{KL}(\Pi_R\|P_0)
 =t_*\E_{\Pi_R}Z-\log M(t_*)-\log Q_*(B_R)
                  =t_*R+O(\log R).
\]
Its mean lies in $[R,R+1]$. Mixing it with $Q_*$ in the unique
proportion giving mean $s$ and using convexity of relative entropy
matches $t_*s-\log M(t_*)$ in the limit. Equality would require
$Q=Q_*$, whose mean is too small.
\end{proof}

Suppose, more specifically, that the maximizing set is a compact smooth
$m$-dimensional submanifold of the $d$-dimensional sphere, and the
correlation deficit is comparable to the $\beta$th power of its normal
distance. Tubular coordinates then give
$V(u)\asymp u^{(d-m)/\beta}$. Exact asymptotics require the corresponding
leading normal-volume limit, but the two-sided comparison already
suffices for the endpoint integrability test. Thus the exponent is
\[
                  \alpha=\frac{d-m}{\beta}-\frac{d+1}{2}. \tag{30}
\]
An isolated sharp corner has $m=0,\beta=1$ and gives
$\alpha=(d-1)/2$, as for ReLU parity. An isolated smooth nondegenerate
quadratic maximum has $m=0,\beta=2$ and gives $\alpha=-1/2$:
the endpoint normalizer then diverges, even though the maximizing
directions themselves form a finite set. Linear $G(U)$ supplies an
explicit example of the latter case.

For ReLU parity, the corner identity gives
\[
 L\Gamma(d+1)=\rho^d\sum_s\int_{v_s^\perp}e^{-h_s(t)/\rho}\,dt.
\]
Indeed its angular gap is $h_s(u)/\rho^2+o(\|u\|)$, and the
volume of a norm ball scales with the $d$th power of its radius.
Formula (27) therefore reduces exactly to $A=\alpha C$ in (10)--(11).
For leaky ReLU with fixed $0\le\lambda<1$,
$\phi_\lambda(t)=(1-\lambda)\relu(t)+\lambda t$, the parity
statistic for even $r\ge2$ is $(1-\lambda)G(U)$: its affine part
has zero parity coefficient. It has the same exponent and moment
thresholds, with the critical tilt rescaled by $(1-\lambda)^{-1}$.

\end{document}
